% Options for packages loaded elsewhere
\PassOptionsToPackage{unicode}{hyperref}
\PassOptionsToPackage{hyphens}{url}
\documentclass[
]{article}
\usepackage{xcolor}
\usepackage[margin=1in]{geometry}
\usepackage{amsmath,amssymb}
\setcounter{secnumdepth}{-\maxdimen} % remove section numbering
\usepackage{iftex}
\ifPDFTeX
  \usepackage[T1]{fontenc}
  \usepackage[utf8]{inputenc}
  \usepackage{textcomp} % provide euro and other symbols
\else % if luatex or xetex
  \usepackage{unicode-math} % this also loads fontspec
  \defaultfontfeatures{Scale=MatchLowercase}
  \defaultfontfeatures[\rmfamily]{Ligatures=TeX,Scale=1}
\fi
\usepackage{lmodern}
\ifPDFTeX\else
  % xetex/luatex font selection
\fi
% Use upquote if available, for straight quotes in verbatim environments
\IfFileExists{upquote.sty}{\usepackage{upquote}}{}
\IfFileExists{microtype.sty}{% use microtype if available
  \usepackage[]{microtype}
  \UseMicrotypeSet[protrusion]{basicmath} % disable protrusion for tt fonts
}{}
\makeatletter
\@ifundefined{KOMAClassName}{% if non-KOMA class
  \IfFileExists{parskip.sty}{%
    \usepackage{parskip}
  }{% else
    \setlength{\parindent}{0pt}
    \setlength{\parskip}{6pt plus 2pt minus 1pt}}
}{% if KOMA class
  \KOMAoptions{parskip=half}}
\makeatother
\usepackage{color}
\usepackage{fancyvrb}
\newcommand{\VerbBar}{|}
\newcommand{\VERB}{\Verb[commandchars=\\\{\}]}
\DefineVerbatimEnvironment{Highlighting}{Verbatim}{commandchars=\\\{\}}
% Add ',fontsize=\small' for more characters per line
\usepackage{framed}
\definecolor{shadecolor}{RGB}{248,248,248}
\newenvironment{Shaded}{\begin{snugshade}}{\end{snugshade}}
\newcommand{\AlertTok}[1]{\textcolor[rgb]{0.94,0.16,0.16}{#1}}
\newcommand{\AnnotationTok}[1]{\textcolor[rgb]{0.56,0.35,0.01}{\textbf{\textit{#1}}}}
\newcommand{\AttributeTok}[1]{\textcolor[rgb]{0.13,0.29,0.53}{#1}}
\newcommand{\BaseNTok}[1]{\textcolor[rgb]{0.00,0.00,0.81}{#1}}
\newcommand{\BuiltInTok}[1]{#1}
\newcommand{\CharTok}[1]{\textcolor[rgb]{0.31,0.60,0.02}{#1}}
\newcommand{\CommentTok}[1]{\textcolor[rgb]{0.56,0.35,0.01}{\textit{#1}}}
\newcommand{\CommentVarTok}[1]{\textcolor[rgb]{0.56,0.35,0.01}{\textbf{\textit{#1}}}}
\newcommand{\ConstantTok}[1]{\textcolor[rgb]{0.56,0.35,0.01}{#1}}
\newcommand{\ControlFlowTok}[1]{\textcolor[rgb]{0.13,0.29,0.53}{\textbf{#1}}}
\newcommand{\DataTypeTok}[1]{\textcolor[rgb]{0.13,0.29,0.53}{#1}}
\newcommand{\DecValTok}[1]{\textcolor[rgb]{0.00,0.00,0.81}{#1}}
\newcommand{\DocumentationTok}[1]{\textcolor[rgb]{0.56,0.35,0.01}{\textbf{\textit{#1}}}}
\newcommand{\ErrorTok}[1]{\textcolor[rgb]{0.64,0.00,0.00}{\textbf{#1}}}
\newcommand{\ExtensionTok}[1]{#1}
\newcommand{\FloatTok}[1]{\textcolor[rgb]{0.00,0.00,0.81}{#1}}
\newcommand{\FunctionTok}[1]{\textcolor[rgb]{0.13,0.29,0.53}{\textbf{#1}}}
\newcommand{\ImportTok}[1]{#1}
\newcommand{\InformationTok}[1]{\textcolor[rgb]{0.56,0.35,0.01}{\textbf{\textit{#1}}}}
\newcommand{\KeywordTok}[1]{\textcolor[rgb]{0.13,0.29,0.53}{\textbf{#1}}}
\newcommand{\NormalTok}[1]{#1}
\newcommand{\OperatorTok}[1]{\textcolor[rgb]{0.81,0.36,0.00}{\textbf{#1}}}
\newcommand{\OtherTok}[1]{\textcolor[rgb]{0.56,0.35,0.01}{#1}}
\newcommand{\PreprocessorTok}[1]{\textcolor[rgb]{0.56,0.35,0.01}{\textit{#1}}}
\newcommand{\RegionMarkerTok}[1]{#1}
\newcommand{\SpecialCharTok}[1]{\textcolor[rgb]{0.81,0.36,0.00}{\textbf{#1}}}
\newcommand{\SpecialStringTok}[1]{\textcolor[rgb]{0.31,0.60,0.02}{#1}}
\newcommand{\StringTok}[1]{\textcolor[rgb]{0.31,0.60,0.02}{#1}}
\newcommand{\VariableTok}[1]{\textcolor[rgb]{0.00,0.00,0.00}{#1}}
\newcommand{\VerbatimStringTok}[1]{\textcolor[rgb]{0.31,0.60,0.02}{#1}}
\newcommand{\WarningTok}[1]{\textcolor[rgb]{0.56,0.35,0.01}{\textbf{\textit{#1}}}}
\usepackage{longtable,booktabs,array}
\usepackage{calc} % for calculating minipage widths
% Correct order of tables after \paragraph or \subparagraph
\usepackage{etoolbox}
\makeatletter
\patchcmd\longtable{\par}{\if@noskipsec\mbox{}\fi\par}{}{}
\makeatother
% Allow footnotes in longtable head/foot
\IfFileExists{footnotehyper.sty}{\usepackage{footnotehyper}}{\usepackage{footnote}}
\makesavenoteenv{longtable}
\usepackage{graphicx}
\makeatletter
\newsavebox\pandoc@box
\newcommand*\pandocbounded[1]{% scales image to fit in text height/width
  \sbox\pandoc@box{#1}%
  \Gscale@div\@tempa{\textheight}{\dimexpr\ht\pandoc@box+\dp\pandoc@box\relax}%
  \Gscale@div\@tempb{\linewidth}{\wd\pandoc@box}%
  \ifdim\@tempb\p@<\@tempa\p@\let\@tempa\@tempb\fi% select the smaller of both
  \ifdim\@tempa\p@<\p@\scalebox{\@tempa}{\usebox\pandoc@box}%
  \else\usebox{\pandoc@box}%
  \fi%
}
% Set default figure placement to htbp
\def\fps@figure{htbp}
\makeatother
\setlength{\emergencystretch}{3em} % prevent overfull lines
\providecommand{\tightlist}{%
  \setlength{\itemsep}{0pt}\setlength{\parskip}{0pt}}
\usepackage{bookmark}
\IfFileExists{xurl.sty}{\usepackage{xurl}}{} % add URL line breaks if available
\urlstyle{same}
\hypersetup{
  pdftitle={Wasps\_Rule},
  pdfauthor={Tanner Lyons},
  hidelinks,
  pdfcreator={LaTeX via pandoc}}

\title{Wasps\_Rule}
\author{Tanner Lyons}
\date{2026-04-27}

\begin{document}
\maketitle

\section{Section 1: Create, Save, and Knit Your First R Markdown
File}\label{section-1-create-save-and-knit-your-first-r-markdown-file}

Author: Jaslyn Davenport
(\href{mailto:jcdavenport@uncg.edu}{\nolinkurl{jcdavenport@uncg.edu}})

Objective: Learn how to create and generate a basic R Markdown document.

\ldots\ldots\ldots\ldots\ldots\ldots\ldots\ldots.

\subsection{1.1 How I Created a New R Markdown
File}\label{how-i-created-a-new-r-markdown-file}

\ldots\ldots\ldots\ldots\ldots\ldots\ldots\ldots.

To begin, open RStudio and create a new R Markdown file.\\
Go to \textbf{File \textgreater{} New File \textgreater{} R Markdown}.

QUESTION: What do you think happens when the new file is opened?\\
ANSWER: A new, blank R markdown file is created.

\ldots\ldots\ldots\ldots\ldots\ldots\ldots\ldots.

\subsection{1.2 Install and Load
Packages}\label{install-and-load-packages}

\ldots\ldots\ldots\ldots\ldots\ldots\ldots\ldots.

Before knitting your file, make sure the required packages are installed
and loaded.

TASK: Run the following in the Console:

install.packages(``rmarkdown'') install.packages(``knitr'')
library(rmarkdown) library(knitr) install.packages(``tinytex'')
tinytex::install\_tinytex()

QUESTION: Why are these packages needed? ANSWER: These packages are used
to format markdown text and combine markdown elements into single,
cohesive outputs.

\ldots\ldots\ldots\ldots\ldots\ldots\ldots\ldots. \#\# 1.3 Modify the
YAML Header

At the top of the file, you will see the YAML header between two sets of
--- This section controls the document settings.

TASK: Change the YAML header so that it includes: a new title your name
as the author the date PDF output

Example: title: ``My First R Markdown File'' author: ``Jaslyn
Davenport'' date: ``April 2026'' output: pdf\_document

QUESTION: What is the purpose of the YAML header? ANSWER: The YAML
header attributes the authors and provides formatting context.

\ldots\ldots\ldots\ldots\ldots\ldots\ldots\ldots. \#\#1.4 Save and Knit
file

After editing the file, save it as an .Rmd file. Then click the Knit
button at the top of RStudio.

If you get an error message about LaTeX or TinyTeX, install TinyTeX by
running: install.packages(``tinytex'') tinytex::install\_tinytex() Then
click Knit again.

QUESTION: What happens when you click Knit? ANSWER:

QUESTION: What type of file is created? ANSWER:

\ldots\ldots\ldots\ldots\ldots\ldots\ldots\ldots. \#\#1.5 Final
observation After knitting, compare the .Rmd file to the final PDF.

QUESTION: How is the output file different from the source file? ANSWER:

\begin{center}\rule{0.5\linewidth}{0.5pt}\end{center}

\section{Section 2: Formatting Text Using RMarkdown
Syntax}\label{section-2-formatting-text-using-rmarkdown-syntax}

\begin{center}\rule{0.5\linewidth}{0.5pt}\end{center}

Author: Arnav Pandey
(\href{mailto:a_pandey2@uncg.edu}{\nolinkurl{a\_pandey2@uncg.edu}})

Objective: learn how to format the appearance of your document in your
chosen output format.

\ldots\ldots\ldots\ldots\ldots\ldots\ldots\ldots.

\subsection{2.1. Headings}\label{headings}

\ldots\ldots\ldots\ldots\ldots\ldots\ldots\ldots.

TASK: Read the question below, click \textbf{Knit} to render the
document and return here to answer the question.

QUESTION: What difference do you notice about the above line that starts
with
\texttt{\#\ Section\ 2:\ Formatting\ a\ Report\ Using\ RMarkdown\ Syntax}
compared to the line
\texttt{Objective:\ learn\ to\ format\ how\ your\ document\ appears\ in\ the\ output\ of\ your\ choice}?\\
ANSWER:

\hfill\break

Placing a hashtag (\#) at the beginning of a line turns it into a
heading. RMarkdown supports six levels of headings.

\hfill\break

TASK: Click Knit to observe how different heading levels below appear in
the rendered output.

\section{Line 1: This is an example to show different levels of headings
in
RMarkdown.}\label{line-1-this-is-an-example-to-show-different-levels-of-headings-in-rmarkdown.}

\subsection{Line 2: Line 1: This is an example to show different levels
of headings in
RMarkdown.}\label{line-2-line-1-this-is-an-example-to-show-different-levels-of-headings-in-rmarkdown.}

\subsubsection{Line 3: This is an example to show different levels of
headings in
RMarkdown.}\label{line-3-this-is-an-example-to-show-different-levels-of-headings-in-rmarkdown.}

\paragraph{Line 4: This is an example to show different levels of
headings in
RMarkdown.}\label{line-4-this-is-an-example-to-show-different-levels-of-headings-in-rmarkdown.}

\subparagraph{Line 5: This is an example to show different levels of
headings in
RMarkdown.}\label{line-5-this-is-an-example-to-show-different-levels-of-headings-in-rmarkdown.}

Line 6: This is an example to show different levels of headings in
RMarkdown.

QUESTION: What happens to the size of text when you increase the number
of hashtags (\#)?\\
ANSWER:

\hfill\break

In the example above, the higher the number of hashtags, the lower the
heading level (smaller font).

QUESTION: How many hashtags would you use if you wanted highest level
heading (largest font)?\\
ANSWER:

\hfill\break

QUESTION: In the rendered output, do Line 5 and Line 6 from above
example appear in the same line?\\
ANSWER:

\hfill\break

TASK: Go to the metadata for this assignment and change the output from
\texttt{pdf\_document} to \texttt{html\_document}.

QUESTION: Did your Line 5 and Line 6 appear in different lines?\\
ANSWER:

\hfill\break

Sometimes the way the program renders output differs between formats.
You can leave two spaces at the end of the line for a line break.

TASK: Change the output back to \texttt{pdf\_document}.

TASK: \textbf{Go back}to all the above example lines ``Line 1'' thru
``Line 6'' within single set of three backticks ```. Go back above and
do as below:

\begin{verbatim}

# Line 1: This is an example to show different levels of headings in RMarkdown.
## Line 2: Line 1: This is an example to show different levels of headings in RMarkdown.
### Line 3: This is an example to show different levels of headings in RMarkdown.
#### Line 4: This is an example to show different levels of headings in RMarkdown.
##### Line 5: This is an example to show different levels of headings in RMarkdown.   
###### Line 6: This is an example to show different levels of headings in RMarkdown.
\end{verbatim}

\ldots\ldots\ldots\ldots\ldots\ldots\ldots\ldots.

\subsection{2.2 Italics}\label{italics}

\ldots\ldots\ldots\ldots\ldots\ldots\ldots\ldots.\\
You may sometimes want to italicize text. RMarkdown allows you to do
this easily. To make your text italics, you can enclose it within single
asterisk \texttt{*{[}text{]}*}.

\emph{Line 1: This is an example to show how to make text italics.}

\hfill\break

TASK: Rewrite Line 1 below and only italicize the word ``italics.''
Leave two spaces at the end of sentence for a line break.\\
ANSWER:

\hfill\break

You may be wondering if there is another way to make the text italics,
and in fact, there is. You can also enclose your text within single
underscore \texttt{\_{[}text{]}\_}.

\hfill\break

TASK: In the paragraph below, make ``Arnav,'' ``Biological Data
Wrangling and Visualization'' and ``RMarkdown'' italics in the following
paragraph using both the methods (asterisks and underscore) mentioned
above at least once.

My name is Arnav Pandey. I am currently enrolled in Biological Data
Wrangling and Visualization (BIO 457/657).This is a TASK: to use
RMarkdown to make a text output italics.

\hfill\break

\ldots\ldots\ldots\ldots\ldots\ldots\ldots\ldots.

\subsection{2.3 Bold}\label{bold}

\ldots\ldots\ldots\ldots\ldots\ldots\ldots\ldots.\\
Sometimes you may want to bold your text to emphasize its importance.
RMarkdown supports that as well.

TASK: Read the question below, click Knit, and return to answer the
question below.

QUESTION: Which of the following lines gives you bold text?\\
Line 1. This is an example to show how to make text bold.\\
\textbf{Line 2: This is an example to show how to make text bold.}\\
ANSWER:

\hfill\break

You can use two asterisks before and after the text to make it bold
(\texttt{**{[}text{]}**}).\\
There is also another way to make a text bold by enclosing your text
within two underscores(\texttt{\_\_{[}text{]}\_\_}).

Don't believe me? Knit to see if the line below is bold.

\textbf{Line 3: This is an example to show how to make text bold.}

\hfill\break

Like you practiced with italics, you can make any word or phrase within
a sentence bold. See example of it below when you hit Knit.\\
Line 5: This is an example to show how to make text \textbf{bold}.\\
Line 4: This is an \textbf{example} to show how to make text bold.

\hfill\break

TASK: In the paragraph below, make ``Arnav,'' ``Biological Data
Wrangling and Visualization'' and ``RMarkdown'' bold in the following
paragraph using both the methods (asterisks and underscores) mentioned
above at least once.\\
My name is Arnav Pandey. I am currently enrolled in Biological Data
Wrangling and Visualization (BIO 457/657). This is a TASK: to use
RMarkdown to make a text output bold.

\hfill\break

\ldots\ldots\ldots\ldots\ldots\ldots\ldots\ldots.

\subsection{2.4 Bold and Italics}\label{bold-and-italics}

\ldots\ldots\ldots\ldots\ldots\ldots\ldots\ldots.

Is there a way to make a text both bold and italics at the same time?
Indeed, there is! You can enclose your text within three asterisks
(\texttt{***{[}text{]}***}). or underscores
(\texttt{\_\_\_{[}text{]}\_\_\_}).

\hfill\break

TASK: In the paragraph below, make ``Arnav,'' ``Biological Data
Wrangling and Visualization'' and ``RMarkdown'' bold italics using both
the methods mentioned above at least once.\\
My name is Arnav Pandey. I am currently enrolled in Biological Data
Wrangling and Visualization (BIO 457/657). This is a TASK: to use
RMarkdown to make a text output bold.

\hfill\break

In sections 2.2, 2.3, and 2.4, you learnt how to make text italics,
bold, or bold+italics by enclosing the text within one, two, and three
asterisks or underscores, respectively. You might be tempted to start
enclosing a text by underscore and end it using asterisk. For example:

*italics\_\\
\_\_bold**\\
***bold italics\_\_\_

Above lines will not give your desired output, so you will have to be
consistent with how you enclose your texts. Use the same symbols at both
ends!

\ldots\ldots\ldots\ldots\ldots\ldots\ldots\ldots.

\subsection{2.5 Bulleted list}\label{bulleted-list}

\ldots\ldots\ldots\ldots\ldots\ldots\ldots\ldots.

In RMarkdown, you can list your items in the RMarkdown as an unordered
(bulleted) list.

To do that, you can follow these steps. Step 1: leave an empty line
before list. Step 2: Use `+', `-', or '*' followed by single space
before you write your list item.\\
Step 3: To add another item, leave two spaces at the end of the previous
item and go to a new line below and repeat Step 2.

TASK: Knit to see how the items below appear in a list in the output.

\begin{itemize}
\tightlist
\item
  Apple
\item
  Banana
\item
  Raspberry
\end{itemize}

You can also nest subitems within an item. If you want a nested list as
such, just indent the subitem that you want within an item with two
spaces. Below is an example of it.

TASK: Hit Knit to look at the output of the lines below:

\begin{itemize}
\tightlist
\item
  Apple
\item
  Banana
\item
  Berries

  \begin{itemize}
  \tightlist
  \item
    Blueberries
  \item
    Strawberries
  \item
    Raspberries
  \item
    Blackberries
  \item
    Cranberries
  \end{itemize}
\item
  Papaya
\end{itemize}

\hfill\break

TASK: Copy the above list below, add an item ``Melon'' and nest
different types of melons as its subitems. Hit Knit to see your output.

\ldots\ldots\ldots\ldots\ldots\ldots\ldots\ldots.

\subsection{2.6 Numbered List}\label{numbered-list}

\ldots\ldots\ldots\ldots\ldots\ldots\ldots\ldots.

Getting a numbered list is even easier compared to unordered or bulleted
list. To generate a numbered list, you can follow these steps. Step 1:
Leave an empty line before list. Step 2: Use uppercase/lowercase English
alphabets, Arabic numerals, uppercase or lowercase roman numerals
followed by space before you write your list item. Be consistent how you
want to list. Step 3: To add another item, leave two spaces at the end
of the previous item and go to a new line below and repeat Step 2.

I. Apple\\
II. Banana\\
III. Raspberry

\hfill\break

You can also nest subitems within an item. If you want a nested list as
such, just indent the subitem that you want within an item with three
spaces. You can use different order format for your subitems than your
primary item as well. Below is an example of it. \textbf{Please be
careful with spaces!}

\begin{enumerate}
\def\labelenumi{\arabic{enumi}.}
\tightlist
\item
  Apple
\item
  Banana
\item
  Berries

  \begin{enumerate}
  \def\labelenumii{\arabic{enumii}.}
  \tightlist
  \item
    Blueberries
  \item
    Strawberries
  \item
    Raspberries
  \item
    Blackberries
  \item
    Cranberries
  \end{enumerate}
\item
  Papaya
\end{enumerate}

Note: If you want to use Roman Numerals, you may have to leave an empty
line between items on list.

\begin{center}\rule{0.5\linewidth}{0.5pt}\end{center}

\section{Section 3: Inserting Code Chunks and Running Basic R
Code}\label{section-3-inserting-code-chunks-and-running-basic-r-code}

\begin{center}\rule{0.5\linewidth}{0.5pt}\end{center}

Author: Michelle Bautista Canuto
(\href{mailto:m_bautistac@uncg.edu}{\nolinkurl{m\_bautistac@uncg.edu}})
We can insert a code chunk by clicking on the green-button with a ``c''
on the inside and a ``+'' on the upper left. Click ``R''. We can insert
functions within the chunk of code. We will see this next.

QUESTION: What is a code chunk? (hint look back at the slides in Canvas)
ANSWER:

TASK: Practice inserting a code chunk. Run a simple calculation such as
``2+2'' below \{r\}. Run this code by clicking on the green
side-triangle button:

Next, let's assign variables. Run the following variables:

\begin{Shaded}
\begin{Highlighting}[]
\NormalTok{x}\OtherTok{\textless{}{-}}\DecValTok{10}
\NormalTok{y}\OtherTok{\textless{}{-}}\DecValTok{5}
\end{Highlighting}
\end{Shaded}

To view the contents of the variable x, just type `x' and run

\begin{Shaded}
\begin{Highlighting}[]
\NormalTok{x}
\end{Highlighting}
\end{Shaded}

\begin{verbatim}
## [1] 10
\end{verbatim}

TASK: Write a chunk of code mulitpliying the variables x and y.

QUESTION: What was the result? ANSWER:

We will analyze the built-in `mtcars' dataset. Run the following chunk
of code.

\begin{Shaded}
\begin{Highlighting}[]
\FunctionTok{head}\NormalTok{(mtcars)}
\end{Highlighting}
\end{Shaded}

\begin{verbatim}
##                    mpg cyl disp  hp drat    wt  qsec vs am gear carb
## Mazda RX4         21.0   6  160 110 3.90 2.620 16.46  0  1    4    4
## Mazda RX4 Wag     21.0   6  160 110 3.90 2.875 17.02  0  1    4    4
## Datsun 710        22.8   4  108  93 3.85 2.320 18.61  1  1    4    1
## Hornet 4 Drive    21.4   6  258 110 3.08 3.215 19.44  1  0    3    1
## Hornet Sportabout 18.7   8  360 175 3.15 3.440 17.02  0  0    3    2
## Valiant           18.1   6  225 105 2.76 3.460 20.22  1  0    3    1
\end{verbatim}

QUESTION: What do you think the head() function does? ANSWER:

To display the summary statistics for the dataset use the summary()
function. Run the following chunk of code.

\begin{Shaded}
\begin{Highlighting}[]
\FunctionTok{summary}\NormalTok{(mtcars)}
\end{Highlighting}
\end{Shaded}

\begin{verbatim}
##       mpg             cyl             disp             hp       
##  Min.   :10.40   Min.   :4.000   Min.   : 71.1   Min.   : 52.0  
##  1st Qu.:15.43   1st Qu.:4.000   1st Qu.:120.8   1st Qu.: 96.5  
##  Median :19.20   Median :6.000   Median :196.3   Median :123.0  
##  Mean   :20.09   Mean   :6.188   Mean   :230.7   Mean   :146.7  
##  3rd Qu.:22.80   3rd Qu.:8.000   3rd Qu.:326.0   3rd Qu.:180.0  
##  Max.   :33.90   Max.   :8.000   Max.   :472.0   Max.   :335.0  
##       drat             wt             qsec             vs        
##  Min.   :2.760   Min.   :1.513   Min.   :14.50   Min.   :0.0000  
##  1st Qu.:3.080   1st Qu.:2.581   1st Qu.:16.89   1st Qu.:0.0000  
##  Median :3.695   Median :3.325   Median :17.71   Median :0.0000  
##  Mean   :3.597   Mean   :3.217   Mean   :17.85   Mean   :0.4375  
##  3rd Qu.:3.920   3rd Qu.:3.610   3rd Qu.:18.90   3rd Qu.:1.0000  
##  Max.   :4.930   Max.   :5.424   Max.   :22.90   Max.   :1.0000  
##        am              gear            carb      
##  Min.   :0.0000   Min.   :3.000   Min.   :1.000  
##  1st Qu.:0.0000   1st Qu.:3.000   1st Qu.:2.000  
##  Median :0.0000   Median :4.000   Median :2.000  
##  Mean   :0.4062   Mean   :3.688   Mean   :2.812  
##  3rd Qu.:1.0000   3rd Qu.:4.000   3rd Qu.:4.000  
##  Max.   :1.0000   Max.   :5.000   Max.   :8.000
\end{verbatim}

Within a new chunk of code, we will code to import the entire mtcars
dataset. Run the following code by clicking on the green side-triangle
button for ``Run Current Chunk''

\begin{Shaded}
\begin{Highlighting}[]
\FunctionTok{data}\NormalTok{(mtcars)}
\end{Highlighting}
\end{Shaded}

QUESTION: After running the previous code, view the dataframe of mtcars.
How many variables are in the mtcars dataset? ANSWER:

Or you can view the dataset in the console by typing ``mtcars'' and
click `Run Selected Lines'. mtcars

To find the number of rows in the dataset, we can use the nrow()
function within a new chunk of code. Run the following code.

\begin{Shaded}
\begin{Highlighting}[]
\FunctionTok{nrow}\NormalTok{(mtcars)}
\end{Highlighting}
\end{Shaded}

\begin{verbatim}
## [1] 32
\end{verbatim}

TASK: Write a code to find the number of columns in the `mtcars' dataset
using the ncol() function. Remember to use a new chunk of code.

QUESTION: What is the result showing in your script? ANSWER:

The `disp' column in the dataset stands for engine displacement. Engine
displacement tells you how big the engine is in cubic inches. The `mpg'
column stands for fuel efficiency in miles per gallon. Lets plot a graph
showing fuel efficiency over displacement for each car. Run the
following code. plot(mtcars\(disp, mtcars\)mpg)

QUESTION: What trend is the plot showing? ANSWER:

We can also use ggplot to run this within our chunk of code. This
permits our plot to be displayed in the final report once ``knit''.

\begin{Shaded}
\begin{Highlighting}[]
\FunctionTok{library}\NormalTok{(ggplot2)}
\FunctionTok{ggplot}\NormalTok{(mtcars, }\FunctionTok{aes}\NormalTok{(}\AttributeTok{x=}\NormalTok{disp, }\AttributeTok{y=}\NormalTok{mpg)) }\SpecialCharTok{+} \FunctionTok{geom\_point}\NormalTok{()}
\end{Highlighting}
\end{Shaded}

\pandocbounded{\includegraphics[keepaspectratio]{12_RMarkdown_files/12_RMarkdown_files/figure-latex/unnamed-chunk-6-1.pdf}}

Let's change the labels of the x-axis to `Engine Size (cubic inches)'
Let's change the y-axis to `Miles Per Gallon' Let's add a title called
``Fuel Efficiency vs Engine Size''.

\begin{Shaded}
\begin{Highlighting}[]
\FunctionTok{ggplot}\NormalTok{(mtcars,}\FunctionTok{aes}\NormalTok{(}\AttributeTok{x=}\NormalTok{disp, }\AttributeTok{y=}\NormalTok{mpg))}\SpecialCharTok{+}\FunctionTok{geom\_point}\NormalTok{()}\SpecialCharTok{+}
  \FunctionTok{labs}\NormalTok{(}
    \AttributeTok{x=}\StringTok{"Engine Size (cubic inches)"}\NormalTok{,}
    \AttributeTok{y=}\StringTok{"Miles Per Gallon"}\NormalTok{,}
    \AttributeTok{title=}\StringTok{"Fuel Efficiency vs Engine Size"}
\NormalTok{  )}
\end{Highlighting}
\end{Shaded}

\pandocbounded{\includegraphics[keepaspectratio]{12_RMarkdown_files/12_RMarkdown_files/figure-latex/unnamed-chunk-7-1.pdf}}

We could add a legend to show which car each point corresponds to. But
first, we should add a column displaying the car names using dplyr.

\begin{Shaded}
\begin{Highlighting}[]
\FunctionTok{library}\NormalTok{(dplyr)}
\end{Highlighting}
\end{Shaded}

\begin{verbatim}
## 
## Attaching package: 'dplyr'
\end{verbatim}

\begin{verbatim}
## The following objects are masked from 'package:stats':
## 
##     filter, lag
\end{verbatim}

\begin{verbatim}
## The following objects are masked from 'package:base':
## 
##     intersect, setdiff, setequal, union
\end{verbatim}

\begin{Shaded}
\begin{Highlighting}[]
\NormalTok{mtcars}\OtherTok{\textless{}{-}}\NormalTok{mtcars }\SpecialCharTok{\%\textgreater{}\%} 
  \FunctionTok{mutate}\NormalTok{(}\AttributeTok{car\_name=}\FunctionTok{rownames}\NormalTok{(mtcars))}
\end{Highlighting}
\end{Shaded}

Lets'bold' the title Let's add a legend to differentiate each point.

\begin{Shaded}
\begin{Highlighting}[]
\FunctionTok{ggplot}\NormalTok{(mtcars,}\FunctionTok{aes}\NormalTok{(}\AttributeTok{x=}\NormalTok{disp,}\AttributeTok{y=}\NormalTok{mpg, }\AttributeTok{color=}\FunctionTok{factor}\NormalTok{(car\_name)))}\SpecialCharTok{+}
  \FunctionTok{geom\_point}\NormalTok{()}\SpecialCharTok{+}
  \FunctionTok{labs}\NormalTok{(}
    \AttributeTok{x=}\StringTok{"Engine Size (cubic inches)"}\NormalTok{,}
    \AttributeTok{y=}\StringTok{"Miles Per Gallon"}\NormalTok{,}
    \AttributeTok{title=}\StringTok{"Fuel Efficiency vs Engine Size"}\NormalTok{,}
         \AttributeTok{color=}\StringTok{"Car Types"}\NormalTok{)}\SpecialCharTok{+}
  \FunctionTok{theme}\NormalTok{(}\AttributeTok{plot.title =} \FunctionTok{element\_text}\NormalTok{(}\AttributeTok{face=}\StringTok{"bold"}\NormalTok{))}
\end{Highlighting}
\end{Shaded}

\pandocbounded{\includegraphics[keepaspectratio]{12_RMarkdown_files/12_RMarkdown_files/figure-latex/unnamed-chunk-9-1.pdf}}

Lets organize the fuel efficiency from highest to lowest. Run the
following.

\begin{Shaded}
\begin{Highlighting}[]
\NormalTok{mtcars }\SpecialCharTok{\%\textgreater{}\%}
  \FunctionTok{select}\NormalTok{(mpg,disp) }\SpecialCharTok{\%\textgreater{}\%} 
  \FunctionTok{arrange}\NormalTok{(}\FunctionTok{desc}\NormalTok{(mpg))}
\end{Highlighting}
\end{Shaded}

\begin{verbatim}
##                      mpg  disp
## Toyota Corolla      33.9  71.1
## Fiat 128            32.4  78.7
## Honda Civic         30.4  75.7
## Lotus Europa        30.4  95.1
## Fiat X1-9           27.3  79.0
## Porsche 914-2       26.0 120.3
## Merc 240D           24.4 146.7
## Datsun 710          22.8 108.0
## Merc 230            22.8 140.8
## Toyota Corona       21.5 120.1
## Hornet 4 Drive      21.4 258.0
## Volvo 142E          21.4 121.0
## Mazda RX4           21.0 160.0
## Mazda RX4 Wag       21.0 160.0
## Ferrari Dino        19.7 145.0
## Merc 280            19.2 167.6
## Pontiac Firebird    19.2 400.0
## Hornet Sportabout   18.7 360.0
## Valiant             18.1 225.0
## Merc 280C           17.8 167.6
## Merc 450SL          17.3 275.8
## Merc 450SE          16.4 275.8
## Ford Pantera L      15.8 351.0
## Dodge Challenger    15.5 318.0
## Merc 450SLC         15.2 275.8
## AMC Javelin         15.2 304.0
## Maserati Bora       15.0 301.0
## Chrysler Imperial   14.7 440.0
## Duster 360          14.3 360.0
## Camaro Z28          13.3 350.0
## Cadillac Fleetwood  10.4 472.0
## Lincoln Continental 10.4 460.0
\end{verbatim}

QUESTION: Which car has the highest fuel efficiency? ANSWER:

QUESTION: What is the size of it's engine? ANSWER:

TASK: Use code to create a plot comparing disp vs hp from the `mtcars'
dataset. With the following:\\
Use the `disp' data for the x-axis and labeled `Engine Size (cubic
inches)' Use the `hp' data for the y-axis and labeled `Horsepower (HP)'
Add an appropriate title for the plot with a `bold' font Add a legend
colored by `car\_name' and labeled ``Car Types''

Lets organize the horsepower from highest to lowest. Run the following.

\begin{Shaded}
\begin{Highlighting}[]
\NormalTok{mtcars }\SpecialCharTok{\%\textgreater{}\%}
  \FunctionTok{select}\NormalTok{(hp,disp) }\SpecialCharTok{\%\textgreater{}\%} 
  \FunctionTok{arrange}\NormalTok{(}\FunctionTok{desc}\NormalTok{(hp))}
\end{Highlighting}
\end{Shaded}

\begin{verbatim}
##                      hp  disp
## Maserati Bora       335 301.0
## Ford Pantera L      264 351.0
## Duster 360          245 360.0
## Camaro Z28          245 350.0
## Chrysler Imperial   230 440.0
## Lincoln Continental 215 460.0
## Cadillac Fleetwood  205 472.0
## Merc 450SE          180 275.8
## Merc 450SL          180 275.8
## Merc 450SLC         180 275.8
## Hornet Sportabout   175 360.0
## Pontiac Firebird    175 400.0
## Ferrari Dino        175 145.0
## Dodge Challenger    150 318.0
## AMC Javelin         150 304.0
## Merc 280            123 167.6
## Merc 280C           123 167.6
## Lotus Europa        113  95.1
## Mazda RX4           110 160.0
## Mazda RX4 Wag       110 160.0
## Hornet 4 Drive      110 258.0
## Volvo 142E          109 121.0
## Valiant             105 225.0
## Toyota Corona        97 120.1
## Merc 230             95 140.8
## Datsun 710           93 108.0
## Porsche 914-2        91 120.3
## Fiat 128             66  78.7
## Fiat X1-9            66  79.0
## Toyota Corolla       65  71.1
## Merc 240D            62 146.7
## Honda Civic          52  75.7
\end{verbatim}

QUESTION: What trend is the plot showing between engine size and
horsepower? ANSWER:

QUESTION: Which cars have the highest and lowest hp? ANSWER:

Don't forget to save and commit your changes.

\begin{longtable}[]{@{}
  >{\raggedright\arraybackslash}p{(\linewidth - 0\tabcolsep) * \real{1.0000}}@{}}
\toprule\noalign{}
\begin{minipage}[b]{\linewidth}\raggedright
\#\# Section 4: Add Visualizations and Tables to Your Report
\end{minipage} \\
\midrule\noalign{}
\endhead
\bottomrule\noalign{}
\endlastfoot
Author: Themesha Parker
(\href{mailto:tlparker@uncg.edu}{\nolinkurl{tlparker@uncg.edu}}) \\
\#\#\# Goal \\
In this section, you will create graphs and a summary table using the
mtcars data set. You will also practice describing patterns that you
observe in the output. \\
Before beginning, make sure you can access the mtcars dataset. Run the
line below to preview the first few rows. \\
head(mtcars) \\
Also make sure you have Tidyverse installed! Run the line of code
below. \\
library(tidyverse) \\
\end{longtable}

\subsection{Part A: Create a
Scatterplot}\label{part-a-create-a-scatterplot}

The variables wt and mpg in the mtcars data set can help us explore the
relationship between car weight and fuel efficiency.

\subsubsection{TASK:}\label{task}

Create a scatter plot with wt on the x-axis and mpg on the y-axis.

\textbf{Hint:} use the plot() function

\begin{Shaded}
\begin{Highlighting}[]
\CommentTok{\# Write your code below this line}
\end{Highlighting}
\end{Shaded}

\subsubsection{TASK:}\label{task-1}

Modify your scatterplot so that it includes: a title, an x-axis label, a
y-axis label, a point color of your choice

\textbf{Hint:} Think back to previous assignments(main=, xlab= ,ylab=,
and col=)

\begin{Shaded}
\begin{Highlighting}[]
\CommentTok{\# Write your updated scatterplot code below this line}
\end{Highlighting}
\end{Shaded}

QUESTION: What does each point in this scatterplot represent? ANSWER:

QUESTION: What relationship do you observe between weight and mpg?
ANSWER:

QUESTION: Does the graph suggest a positive relationship, a negative
relationship, or no relationship? Explain your answer. ANSWER:

\begin{center}\rule{0.5\linewidth}{0.5pt}\end{center}

\subsection{Part B: Create a Boxplot}\label{part-b-create-a-boxplot}

The variable cyl shows the number of cylinders in each car. A boxplot
can help compare fuel efficiency across cylinder groups.

\subsubsection{TASK:}\label{task-2}

Create a boxplot of mpg by cyl using the boxplot function

\begin{Shaded}
\begin{Highlighting}[]
\CommentTok{\# Write your code below this line}
\end{Highlighting}
\end{Shaded}

\subsubsection{TASK:}\label{task-3}

Modify your boxplot so that it includes: a title, an x-axis label, a
y-axis label, and a fill color of your choice

\begin{Shaded}
\begin{Highlighting}[]
\CommentTok{\# Write your updated boxplot code below this line}
\end{Highlighting}
\end{Shaded}

QUESTION: What does each box in this plot represent? ANSWER:

QUESTION: Which cylinder group appears to have the highest fuel
efficiency? ANSWER:

QUESTION: Which cylinder group appears to have the lowest fuel
efficiency? ANSWER:

QUESTION: Based on the boxplot, how does fuel efficiency seem to change
as the number of cylinders increases? ANSWER:

\begin{center}\rule{0.5\linewidth}{0.5pt}\end{center}

\subsection{Part C: Create a Summary
Table}\label{part-c-create-a-summary-table}

You will now create a table showing the average miles per gallon for
each cylinder group.

\subsubsection{TASK:}\label{task-4}

Use dplyr to calculate the average mpg for each value of cyl. Rename the
data avg\_mpg\_by\_cyl

\textbf{Hint:} you need to group the data by cyl then summarize the mean
of mpg.

Use the space below to write your code:

\begin{Shaded}
\begin{Highlighting}[]
\CommentTok{\# Write your code below this line}
\end{Highlighting}
\end{Shaded}

\subsubsection{TASK:}\label{task-5}

Display your results as a clean table.

Use the space below to write your code:

\begin{Shaded}
\begin{Highlighting}[]
\CommentTok{\# Write your code below this line}
\end{Highlighting}
\end{Shaded}

QUESTION: What does this table show? ANSWER:

QUESTION: Which cylinder group has the highest average mpg? ANSWER:

QUESTION: Which cylinder group has the lowest average mpg? ANSWER:

\begin{center}\rule{0.5\linewidth}{0.5pt}\end{center}

\subsection{Part D: Compare the
Visuals}\label{part-d-compare-the-visuals}

Now that you have created both graphs and the table, compare the
information they provide.

QUESTION: How does the scatterplot help you understand the relationship
between weight and fuel efficiency? ANSWER:

QUESTION: How does the boxplot help you compare different groups of
cars? ANSWER:

\begin{center}\rule{0.5\linewidth}{0.5pt}\end{center}

\subsection{Part E: Now you try}\label{part-e-now-you-try}

\begin{enumerate}
\def\labelenumi{\arabic{enumi}.}
\tightlist
\item
  Use a different variable from mtcars to make an additional scatter
  plot.
\item
  Change the color of your scatterplot.
\end{enumerate}

Write your code below:

\begin{Shaded}
\begin{Highlighting}[]
\CommentTok{\# Write your code below this line}
\end{Highlighting}
\end{Shaded}

QUESTION: What did you change, and how did it affect the output? ANSWER:

\begin{center}\rule{0.5\linewidth}{0.5pt}\end{center}

\section{----------------------------------------------------------}\label{section}

\section{Section 5: Improve and Finalize Your R Markdown
Report}\label{section-5-improve-and-finalize-your-r-markdown-report}

\begin{center}\rule{0.5\linewidth}{0.5pt}\end{center}

\textbf{Objective:} Learn how to make your report more polished,
readable, and professional. In this section, we will take everything we
have learned and improve the overall quality of our R Markdown report so
it is easier to read and more useful.

Author: Dhanush Devanand
(\href{mailto:s_devanand@uncg.edu}{\nolinkurl{s\_devanand@uncg.edu}})

\begin{center}\rule{0.5\linewidth}{0.5pt}\end{center}

\subsection{Step 1: Understanding Inline R
Code}\label{step-1-understanding-inline-r-code}

So far, we have been writing code in chunks. However, sometimes we want
to include results directly inside our text. \textbf{Inline R code}
allows us to insert calculations into sentences without a separate code
chunk.

\subsubsection{TASK:}\label{task-6}

Copy and paste the following lines directly into your R Markdown
document, \textbf{outside} of any code chunk, and then Knit:

\begin{verbatim}
The mtcars dataset has 32 rows and 12 columns.
The average miles per gallon (mpg) is 20.09.
\end{verbatim}

\subsubsection{Check}\label{check}

Did the numbers appear automatically in your rendered sentence?

QUESTION: What is inline R code, and how is it different from a regular
code chunk? ANSWER:

QUESTION: Why is inline R code useful when working with a dataset that
might be updated or changed over time? ANSWER:

\subsubsection{TASK:}\label{task-7}

Now try writing your own inline R code sentence. Using the mtcars
dataset, write a sentence that includes the maximum horsepower value
using inline R code. The function \texttt{max()} will be helpful here.

\textbf{Your sentence:} {[}Write your sentence with inline R code
here.{]}

\subsubsection{Why This Matters}\label{why-this-matters}

Inline R code makes your report dynamic. If the dataset changes, your
results will automatically update without you needing to rewrite
anything. This is especially important in professional or scientific
reports where accuracy matters.

\begin{center}\rule{0.5\linewidth}{0.5pt}\end{center}

\subsection{Step 2: Improving Readability by Hiding
Code}\label{step-2-improving-readability-by-hiding-code}

Sometimes, showing all code can make a report look cluttered. If we only
care about the output --- like a graph --- we can hide the code using
the \texttt{echo=FALSE} chunk option.

\subsubsection{TASK:}\label{task-8}

The chunk below currently shows the code when knitted. Add
\texttt{echo=FALSE} inside the curly braces so that only the plot
appears in the output, not the code. Then Knit your document.

\begin{Shaded}
\begin{Highlighting}[]
\FunctionTok{plot}\NormalTok{(mtcars}\SpecialCharTok{$}\NormalTok{wt, mtcars}\SpecialCharTok{$}\NormalTok{mpg,}
     \AttributeTok{main =} \StringTok{"Weight vs MPG"}\NormalTok{,}
     \AttributeTok{xlab =} \StringTok{"Weight (1000 lbs)"}\NormalTok{,}
     \AttributeTok{ylab =} \StringTok{"Miles Per Gallon"}\NormalTok{)}
\end{Highlighting}
\end{Shaded}

\pandocbounded{\includegraphics[keepaspectratio]{12_RMarkdown_files/12_RMarkdown_files/figure-latex/unnamed-chunk-12-1.pdf}}

\textbf{Hint:} Your chunk header should look like this when you are
done:
\texttt{\textasciigrave{}\textasciigrave{}\textasciigrave{}\{r\ echo=FALSE\}}

\subsubsection{Check}\label{check-1}

Do you see the plot in the output but \textbf{not} the code that created
it?

QUESTION: Why might you want to hide code in a final report? ANSWER:

QUESTION: Can you think of a situation where you would want to
\textbf{show} your code in a report rather than hide it? Who might be
reading that report? ANSWER:

\subsubsection{TASK:}\label{task-9}

Now apply \texttt{echo=FALSE} to the ggplot chunk you created back in
Section 3. Knit the document and confirm only the plot appears.

\begin{center}\rule{0.5\linewidth}{0.5pt}\end{center}

\subsection{Step 3: Removing Warnings and
Messages}\label{step-3-removing-warnings-and-messages}

When you load packages or run certain functions, R sometimes prints
extra warnings and messages that are not important to the reader. These
can make your report look messy. You can suppress them using the chunk
options \texttt{message=FALSE} and \texttt{warning=FALSE}.

\subsubsection{TASK: --- Part A: See the messages
first}\label{task-part-a-see-the-messages-first}

Run the chunk below \textbf{without} any chunk options and Knit. Read
through the output carefully and write down any extra messages or
warnings you notice below.

\begin{Shaded}
\begin{Highlighting}[]
\FunctionTok{library}\NormalTok{(tidyverse)}
\end{Highlighting}
\end{Shaded}

\begin{verbatim}
## -- Attaching core tidyverse packages ------------------------ tidyverse 2.0.0 --
## v forcats   1.0.1     v stringr   1.6.0
## v lubridate 1.9.4     v tibble    3.3.1
## v purrr     1.2.1     v tidyr     1.3.2
## v readr     2.1.6     
## -- Conflicts ------------------------------------------ tidyverse_conflicts() --
## x dplyr::filter() masks stats::filter()
## x dplyr::lag()    masks stats::lag()
## i Use the conflicted package (<http://conflicted.r-lib.org/>) to force all conflicts to become errors
\end{verbatim}

\begin{Shaded}
\begin{Highlighting}[]
\FunctionTok{summary}\NormalTok{(mtcars)}
\end{Highlighting}
\end{Shaded}

\begin{verbatim}
##       mpg             cyl             disp             hp       
##  Min.   :10.40   Min.   :4.000   Min.   : 71.1   Min.   : 52.0  
##  1st Qu.:15.43   1st Qu.:4.000   1st Qu.:120.8   1st Qu.: 96.5  
##  Median :19.20   Median :6.000   Median :196.3   Median :123.0  
##  Mean   :20.09   Mean   :6.188   Mean   :230.7   Mean   :146.7  
##  3rd Qu.:22.80   3rd Qu.:8.000   3rd Qu.:326.0   3rd Qu.:180.0  
##  Max.   :33.90   Max.   :8.000   Max.   :472.0   Max.   :335.0  
##       drat             wt             qsec             vs        
##  Min.   :2.760   Min.   :1.513   Min.   :14.50   Min.   :0.0000  
##  1st Qu.:3.080   1st Qu.:2.581   1st Qu.:16.89   1st Qu.:0.0000  
##  Median :3.695   Median :3.325   Median :17.71   Median :0.0000  
##  Mean   :3.597   Mean   :3.217   Mean   :17.85   Mean   :0.4375  
##  3rd Qu.:3.920   3rd Qu.:3.610   3rd Qu.:18.90   3rd Qu.:1.0000  
##  Max.   :4.930   Max.   :5.424   Max.   :22.90   Max.   :1.0000  
##        am              gear            carb         car_name        
##  Min.   :0.0000   Min.   :3.000   Min.   :1.000   Length:32         
##  1st Qu.:0.0000   1st Qu.:3.000   1st Qu.:2.000   Class :character  
##  Median :0.0000   Median :4.000   Median :2.000   Mode  :character  
##  Mean   :0.4062   Mean   :3.688   Mean   :2.812                     
##  3rd Qu.:1.0000   3rd Qu.:4.000   3rd Qu.:4.000                     
##  Max.   :1.0000   Max.   :5.000   Max.   :8.000
\end{verbatim}

\textbf{What messages or warnings did you see?} {[}Write your response
here.{]}

\subsubsection{TASK: --- Part B: Suppress the
messages}\label{task-part-b-suppress-the-messages}

Now add \texttt{message=FALSE} and \texttt{warning=FALSE} to the chunk
header above and Knit again.

\textbf{Hint:} Your chunk header should look like this when you are
done:
\texttt{\textasciigrave{}\textasciigrave{}\textasciigrave{}\{r\ message=FALSE,\ warning=FALSE\}}

\subsubsection{Check}\label{check-2}

Did the extra messages disappear from your rendered output?

QUESTION: What is the difference between \texttt{message=FALSE} and
\texttt{warning=FALSE}? When might you need to use both at the same
time? ANSWER:

QUESTION: Is it always a good idea to suppress warnings? Can you think
of a case where seeing a warning might actually be important? ANSWER:

\subsubsection{Why This Matters}\label{why-this-matters-1}

Hiding unnecessary messages keeps your report clean and professional,
especially when sharing results with someone who does not need to see
behind-the-scenes R output. However, always make sure you understand
what a warning means before choosing to hide it.

\begin{center}\rule{0.5\linewidth}{0.5pt}\end{center}

\subsection{\texorpdfstring{Step 4: Hiding Code AND Output with
\texttt{include=FALSE}}{Step 4: Hiding Code AND Output with include=FALSE}}\label{step-4-hiding-code-and-output-with-includefalse}

In Step 2, you learned that \texttt{echo=FALSE} hides the code but still
shows the output. Sometimes you may want to run a chunk silently ---
hiding both the code and the output entirely. This is done with
\texttt{include=FALSE}.

This is useful for setup steps like loading libraries or defining
variables that need to run but do not need to appear in the report.

\subsubsection{TASK:}\label{task-10}

Run the chunk below as-is and Knit. Note what appears in the output.
Then change the chunk option from \texttt{echo=FALSE} to
\texttt{include=FALSE} and Knit again.

\begin{verbatim}
##   mean_mpg max_hp min_wt
## 1    20.09    335  1.513
\end{verbatim}

\textbf{Hint:} Your chunk header should look like this when you are
done:
\texttt{\textasciigrave{}\textasciigrave{}\textasciigrave{}\{r\ include=FALSE\}}

\subsubsection{Check}\label{check-3}

When using \texttt{include=FALSE}, does anything from this chunk appear
in the rendered output?

QUESTION: What is the difference between \texttt{echo=FALSE} and
\texttt{include=FALSE}? ANSWER:

QUESTION: Give an example of when you would use \texttt{include=FALSE}
in a real report. ANSWER:

\subsubsection{TASK:}\label{task-11}

Now use the \texttt{my\_summary} object created in the chunk above
inside an inline R code sentence below. Even though the chunk is hidden,
the object is still available to use. Copy and paste the lines below
\textbf{outside} of a code chunk and Knit to confirm the values appear.

\begin{verbatim}
The average miles per gallon across all cars is 20.09 and the most powerful engine produces 335 horsepower.
\end{verbatim}

\subsubsection{Check}\label{check-4}

Did the inline values populate correctly even though the chunk that
created them was hidden?

\begin{center}\rule{0.5\linewidth}{0.5pt}\end{center}

\subsection{Step 5: Controlling Figure
Size}\label{step-5-controlling-figure-size}

When creating plots in R Markdown, you can control the size of your
figures directly in the chunk header using \texttt{fig.width} and
\texttt{fig.height}. This helps make your report look clean and
consistent.

\subsubsection{TASK:}\label{task-12}

Run the chunk below and Knit to see the default figure size.

\begin{Shaded}
\begin{Highlighting}[]
\FunctionTok{ggplot}\NormalTok{(mtcars, }\FunctionTok{aes}\NormalTok{(}\AttributeTok{x =}\NormalTok{ wt, }\AttributeTok{y =}\NormalTok{ mpg)) }\SpecialCharTok{+}
  \FunctionTok{geom\_point}\NormalTok{(}\AttributeTok{color =} \StringTok{"steelblue"}\NormalTok{) }\SpecialCharTok{+}
  \FunctionTok{labs}\NormalTok{(}\AttributeTok{title =} \StringTok{"Weight vs Fuel Efficiency"}\NormalTok{,}
       \AttributeTok{x =} \StringTok{"Weight (1000 lbs)"}\NormalTok{,}
       \AttributeTok{y =} \StringTok{"Miles Per Gallon"}\NormalTok{) }\SpecialCharTok{+}
  \FunctionTok{theme\_bw}\NormalTok{()}
\end{Highlighting}
\end{Shaded}

\pandocbounded{\includegraphics[keepaspectratio]{12_RMarkdown_files/12_RMarkdown_files/figure-latex/unnamed-chunk-15-1.pdf}}

\subsubsection{TASK:}\label{task-13}

Now add \texttt{fig.width=8} and \texttt{fig.height=4} to the chunk
header and Knit again.

\textbf{Hint:} Your chunk header should look like this when you are
done:
\texttt{\textasciigrave{}\textasciigrave{}\textasciigrave{}\{r\ fig.width=8,\ fig.height=4\}}

\subsubsection{Check}\label{check-5}

Did the shape of the plot change in your rendered output?

QUESTION: How did changing \texttt{fig.width} and \texttt{fig.height}
affect the appearance of your plot? ANSWER:

QUESTION: Why might consistent figure sizing matter in a professional or
academic report? ANSWER:

\subsubsection{TASK:}\label{task-14}

Try setting \texttt{fig.width=4} and \texttt{fig.height=6} on the same
chunk. Knit and observe how the plot changes. Describe what you notice
below.

\textbf{Observation:} {[}Write your response here.{]}

\begin{center}\rule{0.5\linewidth}{0.5pt}\end{center}

\subsection{Step 6: Adding a Table of
Contents}\label{step-6-adding-a-table-of-contents}

A table of contents makes it easier to navigate your report, especially
as it gets longer. You can add one by modifying the YAML header at the
very top of your \texttt{.Rmd} file.

\subsubsection{TASK:}\label{task-15}

Make sure the YAML header at the very top of your document looks like
this, then Knit:

\begin{Shaded}
\begin{Highlighting}[]
\PreprocessorTok{{-}{-}{-}}
\FunctionTok{title}\KeywordTok{:}\AttributeTok{ }\StringTok{"Group5\_RMarkdown"}
\FunctionTok{author}\KeywordTok{:}\AttributeTok{ }\StringTok{"Arnav Pandey, Themesha Parker, Michelle Bautista Canuto, Dhanush Devanand"}
\FunctionTok{date}\KeywordTok{:}\AttributeTok{ }\StringTok{"2026{-}04{-}26"}
\FunctionTok{output}\KeywordTok{:}
\AttributeTok{  }\FunctionTok{html\_document}\KeywordTok{:}
\AttributeTok{    }\FunctionTok{toc}\KeywordTok{:}\AttributeTok{ }\CharTok{true}
\PreprocessorTok{{-}{-}{-}}
\end{Highlighting}
\end{Shaded}

\subsubsection{Check}\label{check-6}

Do you now see a table of contents at the top of your rendered HTML
document?

QUESTION: How does a table of contents improve the readability of a long
report? ANSWER:

QUESTION: The table of contents is automatically generated from your
headings. What does this tell you about the importance of using proper
heading levels (\texttt{\#}, \texttt{\#\#}, \texttt{\#\#\#}) throughout
your document? ANSWER:

\begin{center}\rule{0.5\linewidth}{0.5pt}\end{center}

\subsection{Step 7: Adding Hyperlinks}\label{step-7-adding-hyperlinks}

You can include clickable links in your report to point readers to
additional resources. The syntax is \texttt{{[}link\ text{]}(URL)}.

\subsubsection{TASK:}\label{task-16}

Add the following clickable link somewhere in your document and Knit:

\texttt{{[}R\ Markdown\ documentation{]}(https://rmarkdown.rstudio.com)}

\subsubsection{Check}\label{check-7}

Does the link appear clickable after knitting?

\subsubsection{TASK:}\label{task-17}

Now add a second hyperlink of your own. Find a webpage related to
something covered in this assignment --- for example, the mtcars
dataset, ggplot2, or R Markdown chunk options --- and add it below using
the same syntax.

\textbf{Your hyperlink:} {[}Write your hyperlink here.{]}

QUESTION: When might adding a hyperlink be useful in a data report?
ANSWER:

\begin{center}\rule{0.5\linewidth}{0.5pt}\end{center}

\subsection{Step 8: Writing a Strong
Conclusion}\label{step-8-writing-a-strong-conclusion}

Every report should end with a short conclusion that summarizes what was
done and what was found.

\subsubsection{TASK:}\label{task-18}

Write a short paragraph (3--5 sentences) below summarizing what you did
in this assignment. Consider:

\begin{itemize}
\tightlist
\item
  What dataset did you use?
\item
  What visualizations or analyses did you create?
\item
  What improvements did you make to the report in Section 5?
\item
  Why is R Markdown useful for communicating data?
\end{itemize}

\textbf{Your conclusion:} {[}Write your response here.{]}

\subsubsection{TASK:}\label{task-19}

Now rewrite at least one sentence from your conclusion so that it
includes inline R code. For example, instead of writing ``the dataset
has 32 rows,'' use \texttt{32} to generate that number automatically.

\textbf{Your updated sentence with inline R code:} {[}Write your
sentence here.{]}

\begin{center}\rule{0.5\linewidth}{0.5pt}\end{center}

\subsection{Step 9: Final Check and
Knit}\label{step-9-final-check-and-knit}

Before submitting, go through the checklist below to make sure
everything is working correctly. In your \texttt{.Rmd} file, replace
each \texttt{{[}\ {]}} with \texttt{{[}x{]}} once you have confirmed
that step.

\begin{itemize}
\tightlist
\item[$\square$]
  Inline R code displays correctly in rendered text
\item[$\square$]
  Plots appear in the output without showing code (\texttt{echo=FALSE})
\item[$\square$]
  A chunk runs silently using \texttt{include=FALSE}
\item[$\square$]
  No unnecessary warnings or messages appear (\texttt{message=FALSE},
  \texttt{warning=FALSE})
\item[$\square$]
  Figure sizes have been adjusted using \texttt{fig.width} and
  \texttt{fig.height}
\item[$\square$]
  Table of contents is visible at the top of the document
\item[$\square$]
  At least two hyperlinks are present and clickable
\item[$\square$]
  Conclusion paragraph is included with at least one inline R code value
\end{itemize}

Once all boxes are checked, click \textbf{Knit} one final time to
generate your completed report.

\end{document}
